% Figure: stress distribution across a curved-beam (crane-hook) section.
% Left panel sketches the hook section with the centroidal axis and the
% shifted neutral axis; right panel plots the hoop stress across the
% section depth, hyperbolic (nonlinear) rather than the straight-beam
% linear line, with the peak at the inner fibre. The neutral axis sits
% inboard of the centroid by the amount e = R_bar - R_n.
\begin{figure}[htbp]
\centering
\begin{tikzpicture}
% Left: section sketch
\begin{scope}[shift={(0,0)}]
  % section rectangle standing for the hook cross-section (inner at bottom)
  \draw[thick, fill=engblue!6] (0,0) rectangle (1.4,4);
  % inner and outer fibre labels
  \node[font=\scriptsize, anchor=east] at (0,0) {inner};
  \node[font=\scriptsize, anchor=east] at (0,4) {outer};
  % centroidal axis
  \draw[engblue, thick, dashed] (-0.4,2.0) -- (1.8,2.0);
  \node[engblue, font=\scriptsize, anchor=west] at (1.85,2.0) {centroid $\bar{R}$};
  % neutral axis, shifted inward (toward inner fibre)
  \draw[engred, thick] (-0.4,1.55) -- (1.8,1.55);
  \node[engred, font=\scriptsize, anchor=west] at (1.85,1.4) {neutral $R_n$};
  % eccentricity bracket
  \draw[engorange, {Latex}-{Latex}] (0.7,1.55) -- (0.7,2.0);
  \node[engorange, font=\scriptsize, anchor=west] at (0.72,1.78) {$e$};
\end{scope}
% Right: stress plot
\begin{scope}[shift={(4.2,0)}]
  \begin{axis}[
    width=7.5cm, height=6cm,
    xlabel={hoop stress $\sigma_\theta$ (MPa)},
    ylabel={position across section},
    ytick=\empty,
    xmin=-120, xmax=220, ymin=0, ymax=4,
    axis lines=left,
    tick label style={font=\scriptsize},
    label style={font=\small},
  ]
  % zero-stress vertical reference
  \draw[enggray, dashed] (axis cs:0,0) -- (axis cs:0,4);
  % hyperbolic curved-beam distribution: tensile peak at inner (bottom),
  % compressive at outer (top), crossing zero at the neutral axis (y=1.55).
  % sigma(y) = k * (1/r - 1/R_n), r running from inner to outer. Rendered
  % from explicit points to keep it safe.
  \addplot[engred, very thick] coordinates {
    (200,0.0) (150,0.4) (108,0.8) (72,1.2) (30,1.55)
    (0,1.55) (-18,2.0) (-45,2.6) (-70,3.2) (-92,4.0)
  };
  \node[engred, font=\scriptsize, anchor=west] at (axis cs:150,0.3) {curved beam};
  % straight-beam linear reference (what My/I would have predicted)
  \addplot[engblue, dashed, thick] coordinates {
    (120,0.0) (0,2.0) (-120,4.0)
  };
  \node[engblue, font=\scriptsize, anchor=east] at (axis cs:-70,3.6) {straight-beam $My/I$};
  % mark neutral axis crossing
  \filldraw[engorange] (axis cs:0,1.55) circle[radius=1.8pt];
  \end{axis}
\end{scope}
\end{tikzpicture}
\caption[Stress across a curved-beam section]{Hoop stress across the
section of a sharply curved member such as a crane hook. The neutral axis
$R_n$ sits inboard of the centroid $\bar{R}$ by the eccentricity $e$,
and the stress varies hyperbolically across the depth rather than
linearly, so the inner fibre carries a stress well above the straight-beam
prediction $My/I$ drawn dashed for comparison. The inner-fibre
over-stress is why a hook cracks at its inside throat and why curved-beam
theory, not the straight-beam formula, sizes hooks, chain links, and
C-frames.}
\label{fig:vol03ch03:curved-beam-stress}
\end{figure}
